% chapters/17/section-05.tex
\section{Rencana Belajar Lanjutan dan Penutup}
\label{sec:rencana-lanjut}

\subsection{Tiga Arah yang Wajar Ditempuh}

Setelah mata kuliah ini, ada tiga arah yang wajar Anda tempuh. Ketiganya
berangkat dari pekerjaan yang sudah Anda buat, dan ketiganya menuntut
kelanjutan yang berbeda.

\begin{description}
  \item[\textbf{Memperdalam satu perkakas.}] Anda memilih Unity dan
        memperdalamnya: animasi berkerangka, pencahayaan, sistem partikel yang
        lebih lengkap, atau penyimpanan kemajuan pemain. Arah ini cocok
        apabila Anda menikmati pekerjaan menyusun dan menyajikan.
  \item[\textbf{Memperdalam algoritma.}] Anda memperdalam bagian yang
        perhitungannya Anda tulis sendiri: pencarian jalur yang lebih baik,
        perilaku kelompok, atau pembangkitan level secara otomatis. Arah ini
        cocok apabila Anda menikmati pekerjaan merancang dan menganalisis.
  \item[\textbf{Membangun karya yang lebih kecil tetapi selesai.}] Anda
        membuat permainan yang ruang lingkupnya jauh lebih kecil daripada
        proyek akhir, tetapi benar-benar selesai dan dapat dimainkan orang
        lain. Arah ini cocok apabila Anda ingin memiliki karya yang dapat
        ditunjukkan.
\end{description}

Ketiganya sah, dan ketiganya tidak saling meniadakan. Yang perlu dihindari
adalah memulai proyek yang jauh lebih besar daripada proyek akhir ini tanpa
menyelesaikan yang lebih kecil terlebih dahulu.

\begin{tipsbox}
Ukuran proyek yang dapat diselesaikan seorang mahasiswa dalam satu semester,
sendiri, biasanya jauh lebih kecil daripada yang ia duga. Sebuah permainan
dengan satu mekanik yang benar-benar selesai lebih berharga daripada sepuluh
mekanik yang setengah jadi.
\end{tipsbox}

\subsection{Menjaga Portofolio}

Simpanlah pekerjaan Anda di tempat yang dapat Anda buka kembali, dan susunlah
sejak sekarang. Tiga hal yang perlu disimpan untuk setiap karya adalah berkas
yang dapat dijalankan, kode sumbernya, dan satu halaman keterangan yang memuat
apa yang dibuat, perkakas yang dipakai, dan apa yang Anda pelajari.

Portofolio itu berguna untuk tiga hal: mengingatkan Anda sendiri atas apa yang
pernah Anda kerjakan, menunjukkan kemampuan Anda kepada orang lain, dan
menjadi titik tolak bagi karya berikutnya. Proyek akhir pada mata kuliah ini
adalah butir pertamanya.

\subsection{Penutup}

Buku ini dimulai dengan pertanyaan apa itu permainan dan bagaimana komputer
menjalankannya, lalu berakhir dengan sebuah prototipe yang diuji oleh orang
lain. Di antaranya ada sebelas bab yang membangun bagian-bagiannya, dan lima
bab yang menggabungkan serta menilainya.

Apabila ada satu hal yang perlu Anda bawa dari seluruh rangkaian itu, ia
bukanlah perintah atau fungsi tertentu. Ia adalah kebiasaan yang muncul
berulang-ulang pada setiap bab: mengukur, bukan menduga. Waktu satu bingkai
diukur, bukan ditebak; jangkauan lompatan dihitung, bukan dikira-kira;
kesulitan level diuji dengan pemeriksaan, bukan dengan perasaan pembuatnya.
Kebiasaan itu berlaku di luar permainan, dan itulah bagian yang paling lama
bertahan setelah Anda lupa nama perintah yang pernah Anda pakai.

\begin{catatabox}
Terima kasih telah menuntaskan seluruh rangkaian bab ini. Apabila Anda
menemukan bagian yang keliru, kurang jelas, atau tidak sesuai dengan
pelaksanaan perkuliahan, sampaikanlah kepada dosen pengampu. Buku ajar
diperbaiki dengan cara itu, sebagaimana permainan diperbaiki dengan
\istilah{playtest}.
\end{catatabox}
